\clearpage
\VerdierPhysicalPageAnchor{0248}%
\begingroup\small\setlength{\parskip}{0pt}\setlength{\emergencystretch}{1.5em}
\setlength{\abovedisplayskip}{4pt plus 2pt minus 2pt}%
\setlength{\belowdisplayskip}{4pt plus 2pt minus 2pt}%

\VerdierSemanticAnchor{ch4d-4-6-8}
\noindent\textbf{命题 4.6.8.}\quad
\textit{设 $Y$ 是 $\mathcal A$ 的一个复形，且
$\varepsilon\colon Y\longrightarrow I^{\bullet\bullet}$ 是 $Y$ 的一个
Cartan--Eilenberg 内射分解。当 $Y$ 是
$\operatorname{comp}^{+}(\mathcal A)$ 的对象，或者当分解
$I^{\bullet\bullet}$ 具有有限长度时，单复形
$\int_2 I^{\bullet\bullet}$ 是一个内射型复形，并且复形态射}

\VerdierSemanticAnchor{ch4d-4-6-8-morphism}
\[
\textit{$\int_2\varepsilon\colon
Y\longrightarrow\int_2 I^{\bullet\bullet}$}
\]
\textit{是拟同构。}

\VerdierSemanticAnchor{ch4d-4-6-8-proof-spectral-objects}
事实上，态射 $\varepsilon$ 诱导了 $Y$ 的、取值于
$D(\mathcal A)$ 的第一谱对象 (4.1.5) 与由按第一次数滤过的
$I^{\bullet\bullet}$ 所定义的、取值于 $D(\mathcal A)$ 的谱对象
((4.1.7), (4.1.8)) 之间的一个态射。典范上同调函子
$H\colon D(\mathcal A)\longrightarrow\mathcal A$ 把这两个谱对象变为取值于
阿贝尔范畴 $\mathcal A$ 的谱对象。只须证明 $\varepsilon$ 诱导这两个谱对象
之间的同构。由 $Y$ 定义的谱对象总是定常的（第二章，4.4.2），而在上述假设下，
由 $I^{\bullet\bullet}$ 定义的谱对象也是定常的。因此只须证明
$\varepsilon$ 在初始项上诱导同构。这些初始项一方面是

\VerdierSemanticAnchor{ch4d-4-6-8-initial-y}
\[
\mathrm I_1^{p,q}=H^q(Y^p),
\]
另一方面是

\VerdierSemanticAnchor{ch4d-4-6-8-initial-i}
\[
{}'\mathrm I_1^{p,q}=H^q(I^{p,\bullet}).
\]
此外，$\varepsilon$ 在初始项上诱导的态射，是由复形态射

\VerdierSemanticAnchor{ch4d-4-6-8-epsilon-p}
\[
\varepsilon^p\colon Y^p\longrightarrow I^{p,\bullet}.
\]
取上同调而得到的。

\VerdierSemanticAnchor{ch4d-4-6-8-proof-end}
但这个态射是拟同构，这就证明了命题。

\VerdierSemanticAnchor{ch4d-4-6-9}
\medskip\noindent\textbf{4.6.9.}\quad
设 $F\colon D(\mathcal A)\longrightarrow\mathcal B$ 是上同调函子，
$Y$ 是 $\operatorname{comp}(\mathcal A)$ 的一个对象，并且
$Y\longrightarrow I^{\bullet\bullet}$ 是一个 Cartan--Eilenberg 内射分解。
函子 $F$ 把相应的 Cartan--Eilenberg 谱对象 (4.6.7) 变为取值于
$\mathcal B$ 的谱对象，从而得到称为\emph{Cartan--Eilenberg 谱序列}的
一个谱序列。以 $(\mathrm{C.E.})_r^{p,q}$ 表示此谱序列的各项
（第二章，4.3.3.1）。此谱序列的初始项为 (4.1.9.1)

\VerdierSemanticAnchor{ch4d-4-6-9-initial-terms}
\[
(\mathrm{C.E.})_1^{p,q}
=F^{p+q}\bigl(I^{\bullet,p}[-p]\bigr)
=F^q(I^{\bullet,p}),
\]
\endgroup
